\scanpage{9}
\deleted{À tout ``cône'' \unsure{str.}}

Si nous nous intéressons aux structures d'effectivité \inserted{qui peuvent se} définir \deleted{\unsure{pr|par}} en termes de \deleted{certaines} une famille de \uline{$i$} comme intersection des $P(i)$, \deleted{\unsure{dans ce}} la donnée d'une telle structure équivaut à celle d'un ``cône'' dans $\Gamma'$, invariant pr $\pi_1$, contenu \inserted{``\uline{strictement}''} dans celui défini pr $\varepsilon$, contenant la classe de $\unsure{i|j}$, satisfaisant à une condition de \deleted{aplatissement} saturation que nous ne précisons pas ici (est-elle automatiquement satisfaite ?).

\margin{\unsure{(*)} i.e. $\varepsilon$ est ``\uline{intérieur}'' au cône polaire !}

Le cas le plus intéressant est celui où $M^+$ \deleted{$\unsure{Q}$} est défini comme $P(i_1)$, où $i_1$ est comme dans $3^\circ$ \deleted{[\unsure{en fait bien}}, \inserted{[N.B.} le cône \uline{invariant} engendré contient bien $j$, et est contenu dans le $\frac12$-espace défini pr $\varepsilon$, puisq.
\[
i_1\bar i_1=j,\qquad \varepsilon i_1=\mathrm{id}
\]
]. On trouve ainsi une application
\diagram{\begin{tikzcd}[ampersand replacement=\&]
\mathrm{I}(G) \arrow[r] \& \deleted{\unsure{\mathrm{Pol}}}\mathrm{Eff}(G)
\end{tikzcd}}

\mspage{8}
